\chapter{Solid State Theory: Midterm Exam, April 26, 2017}
\label{q:theory-2017-04-midterm}
\begin{sourceNote}
Reconstructed from two printed question pages; the marks total 100\%. Handwritten answers accompanying the source are not part of the questions and are not reproduced. Vectors are consistently typeset in bold.
\end{sourceNote}

\section*{1. (10\%) True or false. Explain for your answer.}
\begin{enumerate}[label=(\alph*)]
\item (5\%) Two atoms of the same type in a crystal must be connected by a translation vector of the underlying Bravais lattice.
\item (5\%) Consider a $\mathbf k$ point lying outside of the first Brillouin zone of a Bravais lattice. There must exist a reciprocal lattice point to which the $\mathbf k$ point is closer than to the origin.
\end{enumerate}

\section*{2. (20\%)}
Consider a Bravais lattice spanned by the primitive vectors
\[
\mathbf t_1=\frac a2(0,1,1),\qquad
\mathbf t_2=\frac a2(1,0,1),\qquad
\mathbf t_3=\frac a2(1,1,0).
\]
\begin{enumerate}[label=(\alph*)]
\item (5\%) Consider a family of direct lattice planes that is normal to $\frac12\mathbf g_1+\mathbf g_2-\mathbf g_3$. Find the distance between two successive lattice planes.
\item (10\%) The lattice points on each lattice plane in (a) form a two-dimensional lattice. Find the density of lattice points (i.e., the number of lattice points per unit area) on the two-dimensional lattice plane.
\item (5\%) Consider a family of direct lattice planes that is normal to $\mathbf t_1+2\mathbf t_2-2\mathbf t_3$. Find the distance between two successive lattice planes.
\end{enumerate}

\section*{3. (25\%)}
Consider the two-dimensional crystal of graphite, which has two carbon atoms per unit cell at positions
\[
\mathbf d_1=0,\qquad \mathbf d_2=a(0,1/\sqrt3).
\]
The primitive translation vectors for the underlying two-dimensional hexagonal lattice are
\[
\mathbf t_1=a(1/2,\sqrt3/2),\qquad
\mathbf t_2=a(-1/2,\sqrt3/2).
\]
Graphite has two $\pi$ bands with dispersion
\[
\epsilon_{\mathbf k}=\epsilon_p\pm\gamma|F_{\mathbf k}|,
\]
where $\epsilon_p$ is a constant, $\gamma>0$, and
\[
|F_{\mathbf k}|^2=1+4\cos(k_xa/2)\cos(\sqrt3k_ya/2)
                   +4\cos^2(k_xa/2).
\]
\begin{enumerate}[label=(\alph*)]
\item (15\%) Identify if or not symmetry points $\Gamma$ ($\mathbf k=0$), $Q$ [$\mathbf k=(2\pi/a)(1/2,1/(2\sqrt3))$], and $P$ [$\mathbf k=(2\pi/a)(2/3,0)$] of the first Brillouin zone are critical points. What type of critical point are they?
\item (10\%) Express the density of states contributed from the states in a small region containing a critical point you choose as you like from those in (a).
\end{enumerate}

\section*{4. (5\%)}
Describe the Bloch theorem.

\section*{5. (5\%)}
Explain the negative value of the exchange energy in a many-electron system from the viewpoint of the difference between the Hartree-Fock and the Hartree theories.

\section*{6. (10\%)}
Describe the Hohenberg-Kohn functional.

\section*{7. (10\%)}
Describe the local density approximation (LDA) in the exchange-correlation functional.

\section*{8. (15\%)}
Write down the Hartree, the Hartree-Fock, and the Kohn-Sham equations and explain each term in the equations. Briefly describe how to solve the equations.
